\documentclass[aps,pre,twocolumn,superscriptaddress,nofootinbib,10pt]{revtex4-2}
\usepackage{graphicx}\usepackage{float}
\usepackage{amsmath,amssymb}
\usepackage{url}\usepackage{xurl}
\newcommand{\lcHtmlSilent}{120}
\newcommand{\lcHtmlRuns}{120}
\newcommand{\lcHtmlFound}{13191}
\newcommand{\lcHtmlTotal}{16876}
\newcommand{\lcHtmlRecall}{0.78}
\newcommand{\lcHtmlRecMin}{0.62}
\newcommand{\lcHtmlRecMax}{0.80}
\newcommand{\wilTwentyLo}{0.84}
\newcommand{\wilFiveLo}{0.57}
\newcommand{\wilZeroTwentyHi}{0.16}
\newcommand{\hsHtmlSilent}{0}
\newcommand{\hsHtmlRuns}{120}
\newcommand{\unHtmlSilent}{0}
\newcommand{\unHtmlRuns}{120}
\newcommand{\mkHtmlSilent}{0}
\newcommand{\mkHtmlRuns}{120}
\newcommand{\dlHtmlSilent}{0}
\newcommand{\dlHtmlRuns}{30}
\newcommand{\liHtmlSilent}{20}
\newcommand{\liHtmlRuns}{120}
\newcommand{\muRotMain}{0.56}
\newcommand{\muRotHold}{0.54}
\newcommand{\muTruncMain}{0.53}
\newcommand{\muTruncHold}{0.52}
\newcommand{\muTruncHoldSilent}{19}
\newcommand{\muTruncHoldRaised}{1}
\newcommand{\pRepaired}{20}
\newcommand{\pFiles}{20}
\newcommand{\pZeroPages}{74}
\newcommand{\pPages}{175}
\newcommand{\cDocxSilent}{0}
\newcommand{\cDocxRuns}{80}
\newcommand{\cPptxSilent}{0}
\newcommand{\cPptxRuns}{80}
\newcommand{\cHtmlSilent}{0}
\newcommand{\cHtmlRuns}{120}
\newcommand{\cTxtSilent}{0}
\newcommand{\cTxtRuns}{100}
\newcommand{\cPdfASilent}{20}
\newcommand{\cPdfARuns}{140}
\newcommand{\cPdfBSilent}{0}
\newcommand{\cPdfBRuns}{140}
\newcommand{\cUnreadPdf}{20}
\newcommand{\glueRealPct}{3.3}
\newcommand{\lcHtmlLossPct}{22}
\newcommand{\pDocxRuns}{80}
\newcommand{\pDocxComplete}{80}
\newcommand{\pDocxIncomplete}{0}
\newcommand{\pDocxRejected}{0}
\newcommand{\pDocxNocore}{0}
\newcommand{\pPptxRuns}{80}
\newcommand{\pPptxComplete}{80}
\newcommand{\pPptxIncomplete}{0}
\newcommand{\pPptxRejected}{0}
\newcommand{\pPptxNocore}{0}
\newcommand{\pHtmlRuns}{120}
\newcommand{\pHtmlComplete}{120}
\newcommand{\pHtmlIncomplete}{0}
\newcommand{\pHtmlRejected}{0}
\newcommand{\pHtmlNocore}{0}
\newcommand{\pTxtRuns}{100}
\newcommand{\pTxtComplete}{40}
\newcommand{\pTxtIncomplete}{0}
\newcommand{\pTxtRejected}{60}
\newcommand{\pTxtNocore}{0}
\newcommand{\pPdfARuns}{120}
\newcommand{\pPdfAComplete}{60}
\newcommand{\pPdfAIncomplete}{0}
\newcommand{\pPdfARejected}{20}
\newcommand{\pPdfANocore}{40}
\newcommand{\pPdfBRuns}{120}
\newcommand{\pPdfBComplete}{60}
\newcommand{\pPdfBIncomplete}{0}
\newcommand{\pPdfBRejected}{20}
\newcommand{\pPdfBNocore}{40}
\newcommand{\eightBitLostLo}{24}
\newcommand{\eightBitLostHi}{26}
\newcommand{\pFourNinety}{106}
\newcommand{\pFourPages}{170}
\newcommand{\pFourDefault}{0.54}
\newcommand{\pFourNoClip}{1.00}
\newcommand{\pFourExtra}{8753}
\newcommand{\pFourWords}{18829}
\newcommand{\preregCells}{74}
\newcommand{\preregMiss}{7}
\newcommand{\preregHit}{67}
\newcommand{\preregCellsHS}{22}
\newcommand{\preregCellsUN}{26}
\newcommand{\preregCellsMK}{26}
\newcommand{\preregAHit}{11}
\newcommand{\preregRaisedHit}{2}
\newcommand{\preregBHit}{49}
\newcommand{\preregUnreadHit}{5}

\newif\ifdraftmode\draftmodefalse

\ifdraftmode\AddToHook{shipout/foreground}{\put(0.5\paperwidth,-0.42in){\makebox[0pt]{\small\sffamily DRAFT, not for distribution}}}\fi

\addtolength{\textheight}{12pt}\addtolength{\topmargin}{-6pt}
\begin{document}\linespread{0.975}\selectfont

\title{Default document loaders in AI frameworks can omit expected text without raising an exception:\\ a short measurement note}

\author{Adam Faust}
\affiliation{MOBLUEHQ, Inc., Long Grove, Illinois, USA}
\date{\ifdraftmode Draft v17 of 26 November 2026. DRAFT, not for distribution.\else November 2026\fi}

\begin{abstract}
Software that lets a language model answer questions about a user's files must first read those files, usually with a one-line ``loader'' call taken from a framework's documentation. We ask a simple question: when such a default call returns normally, has it returned all of the text? We built small files in which every word is known (random tokens in PDF, DOCX, PPTX, HTML and plain text), passed them through selected default calls in six frameworks at pinned versions, and counted the known words that came back. The main finding is that each of the six frameworks returned without an exception while lacking known words on at least one tested path (Table~I). The clearest case is LangChain's \texttt{BSHTMLLoader}: in all \lcHtmlRuns\ HTML fixture runs its default merged words across adjacent block boundaries, leaving \lcHtmlLossPct\% of the expected tokens absent as standalone tokens, and no exception was raised. Because our HTML fixtures omit white space between blocks, as minified markup does, we also checked 58 public web pages: \glueRealPct\% of the tokens extracted with a space separator were absent under the default call. Earlier studies score the quality of PDF extraction or describe what Office extraction cannot reach. This note adds a record, call by call and across five formats, of what is returned and what the caller is told, with alternative calls that, in standalone tests on the specified fixtures, recovered the tokens in several classes or rejected the input. Rates are for constructed files and pinned versions; on real documents we measured only disagreement between tools. Maintainers of LangChain, LlamaIndex, Docling, Haystack, Unstructured and MarkItDown were notified on 25 November 2026, with the measurements and exact calls; the test-file bundle accompanies this note.
\end{abstract}

\maketitle

\section{Introduction}
A retrieval or agent system can only answer from the text it was given. That text usually arrives through a \emph{loader}: a short call, often copied from a framework's documentation, that turns a PDF, an Office file, a web page or a text file into a string. If the loader returns only part of a document, or skips a file, and does so without raising an exception, a caller that relies on exceptions has no reason to suspect that downstream answers rest on incomplete input.

\emph{The question.} For named loader calls used with their default arguments, which inputs lose known content with no exception? We answer it on constructed files, where the truth is known exactly: every word in a test file is a random string, so counting the words that come back measures what was lost.

\emph{The finding.} Each of the six frameworks we tested (LangChain, LlamaIndex, Docling, Haystack, Unstructured and MarkItDown) returned normally while lacking known words on at least one tested path (Table~I). The clearest case is in HTML. LangChain's \texttt{BSHTMLLoader} joins the text of adjacent blocks with nothing between them, so the last word of one paragraph and the first word of the next come back as a single string. This happened in all \lcHtmlRuns\ HTML fixture runs and left \lcHtmlLossPct\% of the expected tokens absent as standalone tokens. Most other losses remove whole parts of a file: the headers, footers and nested tables of DOCX files and the speaker notes of PPTX files. Three kinds alter words instead: character entities left undecoded in HTML, words with a non-ASCII character dropped or replaced in text files that are not UTF-8, and the words of rotated PDF pages, returned with their characters reversed, spaced apart or split rather than as the expected tokens. In standalone tests on the specified fixtures, a single argument or a single check recovered the tokens in several classes or turned the silent loss into an explicit rejection (Sec.~IV).

\emph{What is new.} We measure which expected tokens are missing when a specific default loader call returns without raising; the metric itself is not new. Bast and Korzen scored missing words per tool for fourteen PDF extractors on 12,098 arXiv articles \cite{bast2017}, and ParseBench (its Sec.~3.3) uses normalised word recall and precision \cite{parsebench}. Zhao, Pandey and Misra (their Appendix S4.3) describe headers, footers and other parts that a DOCX body walk cannot reach \cite{zhao2026}, and Liu \emph{et al.}\ and Li measure gaps between what Office shows and what extraction returns \cite{liu2026,li2026}. Our DOCX rows repeat the known header and footer omissions for named calls. Across five formats we report token recall and exception outcomes for six pinned frameworks, including HTML block-boundary merges and plain-text decoding losses, and for selected calls we also record other caller-visible signals.

This is a measurement on test files, not a benchmark of extraction quality. The calls are selected examples, the shortest default calls we know for each loader; we did not survey deployed code and make no claim about how often they are used.

\section{Method}
A generator writes one file per class per seed; the classes are listed in Table~S1, and each format has a plain control. Every word is a random string of 8 to 11 ASCII lower-case letters and digits, unique within the seed, so the truth for each file is a known multiset of tokens. The one exception is three of the five text classes (Latin-1, cp1252 and a mix of UTF-8 and Latin-1), where every fourth word has one interior character replaced by a non-ASCII one.

We lower-case the loader output, split it into tokens, and report \emph{token recall}: the share of truth tokens present, counted with multiplicity. Two words joined into one, one word broken in two, or a word whose non-ASCII character was dropped or replaced count as missing. A run is one loader call on one file, and a run has \emph{core loss} if recall is below 1 on the \emph{core content} of its class: the text a reader of the file would expect back, such as page text, body paragraphs, table cells, headers, footers and speaker notes. Text that these tools cannot be expected to return (words that exist only as an image, form field values, words written by a script, text hidden by CSS) is scored separately and is not called a loss here. Only the returned text is scored, not metadata, and full token recall does not establish order, structure or the absence of extra text. The definition for each class, the scored field of each framework and the exact calls are given in Sec.~S11.

\emph{Silent} means that the call returned without raising an exception; that was the systematic measure. Printed or logged messages were not scored in the framework runs, and a few other signals (logs, status fields, the number of documents returned) were checked on selected calls, as the text says.

Each framework default was run on 20 seeds, except Docling, which was run on the first 5 of them for running time. The library calls that sit underneath the frameworks were run on two independent sets of 20 seeds, a main run and a hold-out run. Outcomes within a class rarely vary. Versions are pinned in Table~S6: langchain-community 0.4.2, llama-index-core 0.14.26, docling 2.152.0, haystack-ai 3.5.0, unstructured 0.27.38 and markitdown 0.1.9. Four of the 30 framework--format pairs were not tested: LangChain on PPTX and plain text, Docling on plain text, and Haystack on PPTX (Sec.~S1), so an absent row in Table~I is not a clean result for those. The calls are the shortest ones for each loader; for LangChain they are \texttt{PyPDFLoader}, \texttt{Docx2txtLoader} and \texttt{BSHTMLLoader}, and for Haystack \texttt{PyPDFToDocument}, \texttt{DOCXToDocument}, \texttt{HTMLToDocument} and \texttt{TextFileToDocument} (all in Sec.~S11). For Haystack, Unstructured and MarkItDown we committed a prediction for each loader and class before the run: of \preregCells\ predictions, \preregMiss\ were wrong, and all of them are listed in Sec.~S3. The stored results and the script that writes every number here are archived at the DOI given under Data and code availability.

\section{Results on constructed files}

\begin{table*}[t]
\caption{Framework defaults on constructed files: rows with core loss or rejection only. Denominators are attempted runs; ``Lost, no exception'' counts runs with token recall below 1 on core content among all runs, and ``Raised'' counts runs that threw. Recall is pooled over the runs that returned. Clean cells are in Table~S2, the corrected rotated-page file in Table~S5 and truncated PDFs (P11, damaged-input handling) in Table~S4; image-only and form PDFs (P7, P9) have no core tokens and are scored separately (Sec.~II). Docling was run on 5 seeds, the others on 20.\label{tab:fw}}
\begin{ruledtabular}
\begin{tabular}{llccc}
Default loader & File class & Lost, no exception & Raised & Recall \\
\colrule
LangChain \texttt{BSHTMLLoader} & HTML, all six classes (H1--H6) & 120/120 & 0 & 0.78 \\
LlamaIndex \texttt{SimpleDirectoryReader} & TXT, UTF-16 with BOM (T2) & 20/20 & 0 & 0.00 \\
 & TXT, Latin-1, cp1252, mixed (T3--T5) & 60/60 & 0 & 0.75 \\
 & HTML, character entities (H5)$^{a}$ & 20/20 & 0 & 0.80 \\
Docling \texttt{DocumentConverter} & DOCX, headers and footers (D3) & 5/5 & 0 & 0.83 \\
 & PPTX, speaker notes (S4) & 5/5 & 0 & 0.57 \\
 & PDF, hidden text over image (P6) & 5/5 & 0 & 0.87 \\
Haystack \texttt{DOCXToDocument} & DOCX, headers and footers (D3) & 20/20 & 0 & 0.82 \\
 & DOCX, table nested in a cell (D4) & 20/20 & 0 & 0.87 \\
Haystack \texttt{TextFileToDocument} & TXT, four non-UTF-8 classes (T2--T5)$^{b}$ & 80/80 & 0 & 0.00 \\
Unstructured \texttt{partition} & PPTX, speaker notes (S4) & 20/20 & 0 & 0.57 \\
 & TXT, Latin-1, cp1252, mixed (T3--T5) & 7/60 & 0 & 0.97 \\
MarkItDown \texttt{convert} & DOCX, headers and footers (D3) & 20/20 & 0 & 0.82 \\
 & TXT, Latin-1, cp1252, mixed (T3--T5) & 7/60 & 0 & 0.97 \\
Unstructured, MarkItDown & PDF, rotated pages, corrected file (P4c) & 20/20, 15/20 & 0 & 0.39, 0.61 \\

\end{tabular}
\end{ruledtabular}
\raggedright\footnotesize $^{a}$ With no \texttt{file\_extractor} registered (the default), the reader returns the raw markup of an \texttt{.html} file, so character entities are not decoded. $^{b}$ Haystack skips the file, writes a warning to the log and returns an empty \texttt{documents} list (zero documents, checked); the call does not raise.
\end{table*}

Table~I lists every framework default that lost core content. We describe the losses by kind, starting with the clearest.

\subsection{Words merged in HTML}
LangChain's \texttt{BSHTMLLoader} calls BeautifulSoup's \texttt{get\_text} with a separator that defaults to the empty string. Text in adjacent block elements is then concatenated: \texttt{...g1h4603dsk</p><p>lb3cew357...} came back as \texttt{g1h4603dsklb3cew357}, and the page title was joined to the first word of the body. This happened in \lcHtmlSilent\ of \lcHtmlRuns\ runs, on all six HTML classes including the plain-paragraph control, with no exception and (in a one-file check) nothing printed. Pooled token recall was \lcHtmlRecall\ (\lcHtmlFound\ of \lcHtmlTotal\ tokens). The call is \texttt{BSHTMLLoader(path).load()}; recall is computed from the returned \texttt{page\_content} against the designated core tokens of the HTML fixtures, and the one-character title and the separately scored content (Sec.~S11) are not in the denominator. The characters are all there; what is lost is the word boundary, which matters wherever text is matched word by word (keyword search, exact matching, table cells).

The remedy is one argument. In both standalone BeautifulSoup runs, \texttt{get\_text(" ")} achieved 1.00 core-token recall in every HTML fixture class, and LangChain exposes the same separator through \texttt{get\_text\_separator=" "}; we did not run that setting through LangChain itself. The other five frameworks' HTML readers had no core loss (0 of 120 runs each, 0 of 30 for Docling), except LlamaIndex, which returns raw markup and so lost tokens on the character-entity class only (H5; Table~I, note $a$).

Two cautions apply. The fixtures omit white space between adjacent block tags, as minified pages do; whitespace between adjacent text nodes prevents this particular merge at that boundary. On public pages the measured difference is accordingly smaller. Of 60 public web pages we compared 58 (two had too little text). At least one token of the separator output was absent from the default output in 52 of them, the default output held fewer than 90\% of the separator output's tokens in 4, and \glueRealPct\% of those tokens were absent overall (1046 of 31585; median page token overlap 0.97). That comparison was not pre-registered, and a separator can split a word styled part-way through, so it measures disagreement between two settings, not verified loss.

\subsection{Text files that are not UTF-8}
LlamaIndex's \texttt{SimpleDirectoryReader} returned text without raising for files that are not UTF-8: nothing usable came back from the UTF-16 files, and only the words without a non-ASCII character (75\%) came back from the 8-bit files. A respondent on the LlamaIndex discussion (\#24444) reproduced this and traced it to the reader's defaults (Sec.~S5). Haystack took the opposite course. With default settings, \texttt{TextFileToDocument} logged ``Could not convert file \ldots\ Skipping it'' and returned an empty \texttt{documents} list for all 80 files in the four non-UTF-8 classes. Skipping and logging a bad input is a reasonable design for batch indexing, where one corrupt file should not stop a job; the point for a caller is that the converter returns zero documents without raising. On Latin-1, cp1252 and mixed text files, Unstructured and MarkItDown both lost about a quarter of the tokens in the same 7 of 60 files (pooled recall 0.97 over the 60), in cells we had predicted would be clean.

\subsection{Parts of Office files}
Several defaults returned the body of an Office file and left out a part of it. On 5 seeds, Docling's default text export did not contain DOCX headers and footers or PPTX speaker notes, with status \texttt{SUCCESS} and no errors (checked afterwards); \texttt{export\_to\_markdown()} gave the same recall (Sec.~S9). Docling may leave headers, footers and notes out of the default export on purpose; we report what the call returned, not a defect. Haystack's DOCX converter returned no headers or footers and, for a table nested in a cell of another table, the four cells of the outer table and none of the six cells of the nested one, in every seed, and printed nothing; the pinned converter builds each cell's text from that cell's own paragraphs only. Unstructured's \texttt{partition} omitted PPTX speaker notes, and MarkItDown lost DOCX headers and footers. LlamaIndex read DOCX and PPTX in full. Haystack has no PPTX row (Sec.~S1).

\subsection{PDF: rotated pages, hidden text and damaged files}
On PDFs whose pages carry rotation entries, Unstructured and MarkItDown lost words with no exception: \texttt{partition} in 20 of 20 files and MarkItDown in 15 of 20 (Table~S5). The loss is page-wide: a page's words are all found or none is found as an exact token. Zero exact-token recall did not mean empty output: the characters of the affected pages were returned, but reversed (MarkItDown, 180 degrees), spaced one character apart (Unstructured, 90 and 270 degrees) or reversed or split (MarkItDown, 270 degrees) (Sec.~S5). On one page rotated by 90 degrees, standalone pdfminer recall was 0.00 with default settings and 1.00 with \texttt{detect\_vertical=True}; Unstructured's active strategy was unrecorded. These numbers are for a corrected file: the first rotated-page file was defective and is withdrawn (Sec.~S5). Docling, on 5 seeds, returned part of the known words for pages with an invisible text layer over an image, again with status \texttt{SUCCESS}.

Truncated PDFs (Table~S4) test something different, the handling of a damaged input, and are scored against the intact original. LangChain, Unstructured and MarkItDown raised on all 20 files; LlamaIndex and Haystack returned zero documents without raising; Docling returned \texttt{PARTIAL\_SUCCESS}, so a caller that checks the status is told.

\subsection{The libraries underneath}
Several of these losses coincide with the engine underneath, so a framework row is a library plus a wrapper. We therefore ran the short library calls directly (Table~S4). Reading only \texttt{doc.paragraphs} with python-docx misses every table cell, nested table, header and footer; reading only text-frame shapes with python-pptx misses table cells, grouped shapes and speaker notes. These are properties of the short call, not library faults: a walk that also visits tables, sections, groups and notes read every class in full, as did docx2txt and MarkItDown's PPTX converter. mammoth's \texttt{extract\_raw\_text} returned the body and tables in full and no header or footer text, in 20 of 20 files in each run, the same mechanism as in \cite{zhao2026}.

For truncated files, cutting a PDF at 50 to 70\% of its bytes made pypdf and pypdfium2 raise on every file. PyMuPDF repaired the file and returned the right number of pages, the later ones empty (\pZeroPages\ of \pPages\ pages with no words across the \pFiles\ main-run files; recall \muTruncMain). All \pFiles\ main-run files returned incomplete text without raising or writing a new message to standard error. All \pRepaired\ set \texttt{Document.is\_repaired} and left messages in MuPDF's internal warning buffer. The short call does not look.

\subsection{Why silence matters}
For retrieval-augmented systems a silently missing page or header could be worse than a loud failure, because the resulting index may look complete: every document is present and searchable, and the absence may stay invisible until a question about the missing text is asked. We did not test retrieval or answers, so this is a possible consequence, not a measured one. Several calls returned zero documents, which a caller that counts documents would see.

\section{Alternative extraction settings and explicit rejection}
None of the remedies is new; each is one argument or one check, and we fixed the list before running it on the 20 framework seeds. For HTML, pass a space as the separator (tested with standalone BeautifulSoup, not through the LangChain wrapper). For DOCX, visit tables (nested ones too) and each section's header and footer, not only \texttt{doc.paragraphs}. For PPTX, recurse into group shapes and read table cells and the notes slide. For text, decode strictly (no \texttt{errors="ignore"} or \texttt{"replace"}), honour a byte-order mark or declared encoding, and raise on invalid bytes. For PDF, use PyMuPDF and raise when \texttt{is\_repaired} is true or the warning buffer is non-empty, or use pypdfium2 after checking the end-of-file marker and \texttt{pikepdf.Pdf.check\_pdf\_syntax()}. Some content is simply unreadable by these tools: image-only pages, form field values and script-rendered HTML have no text layer (they need OCR, a form reader or a browser), and an empty string from them is not a blank page.

In Table~S7, the tested standalone alternatives recovered all core tokens on every returned run with core content on the specified fixtures; some of them reject inputs rather than return text. The PDF runs with complete recall were from the plain, blank-page and hidden-text classes; the truncated files were rejected, and the rotated-page file was not tested with these alternatives. Under its encoding-selection rule, the tested text decoder rejected the 60 Latin-1, cp1252 and mixed-encoding fixtures, turning a silent loss into an explicit one. Image-only and form PDFs and script-written HTML words returned nothing (Sec.~S10).

\emph{Real documents.} On public files no truth is known, and we measured only disagreement between tools: against a pypdfium2 reference, LangChain's and LlamaIndex's PDF readers each returned fewer than 90\% of the reference tokens in 16 of 103 PDFs. This is divergence, not loss, and unchecked by hand (Sec.~S6). No DOCX or PPTX real files were tested.

\section{Limitations}
\emph{Constructed files.} The results show that tested defaults can miss expected tokens on the constructed inputs; they do not say how often real corpora contain that structure. Tokens are random ASCII strings (apart from the 8-bit classes), in one font, from one generator per format, on one machine. There is no study of effects on question answering and no human labelling of real documents.

\emph{Versions and scope.} Package versions were pinned within each environment (Table~S6), default arguments only, six frameworks. Haystack, Unstructured and MarkItDown ran in a separate environment, scored per document, not per page. Coverage is in Sec.~S1. Whether a behaviour is a bug, a documented default or a design choice is for the maintainers to say.

\emph{A run-time check comparing two PDF readers.} We tested a check that compares a loader's tokens with those of a second, independent PDF reader at run time. On the constructed PDF pages both readers missed the same visible content, so their agreement did not establish complete extraction (Sec.~S7). The result is scoped to the PDF setup tested.

\section{Disclosure, provenance and author statement}
\emph{Disclosure.} The maintainers of LangChain, LlamaIndex, Docling, Haystack, Unstructured and MarkItDown were notified on 25 November 2026, with the measurements and exact calls; the test-file bundle accompanies this note. We used blueMonster, a tool developed by MOBLUEHQ, Inc., to generate research topics and brainstorm ideas. The harness, the pre-registered plans, the analysis and this draft were written with Claude large language models (Anthropic) under the author's direction. Two further large language models (ChatGPT, OpenAI; Gemini, Google) criticised earlier drafts of this note and shaped the withdrawal of the two-reader check as a general check, the split between caller settings and engine behaviour, and the plain-settings baseline; these were AI-assisted critiques, not peer review, and no human referee has read this draft. A script writes every number from the committed result files. The author takes responsibility for the content.

The author is an amateur. He has no formal training in statistics, econometrics or machine-learning evaluation and has published no research in them. Much of this work was done while improving MOBLUEHQ's software tools, by using them on a real problem, and in order to learn how such research is done. His inexperience should invite heavy scepticism towards the methods and the claims: errors of a kind a trained researcher would not make are possible, and the checks described here lower that risk without removing it. He welcomes corrections.

\emph{Competing interests.} The author is the founder of MOBLUEHQ, Inc., which develops and intends to sell blueMonster.

\section*{Data and code availability}
The generator, the loader calls, the harness, pinned requirements, the stored results and the script that writes every number here are held in the code-and-fixtures bundle (repository: \url{https://github.com/mobluehq-llc/framework-loader-silent-loss}; DOI \url{https://doi.org/10.5281/zenodo.24522148}). The currently public release of that bundle regenerates an earlier version (v9) of this note, not every number here; the files added since, including the two-reader and Unstructured filename-only results of Secs.~S7 and S11, go into its next release. The public real-document files are pinned by checksum and not redistributed.

\begin{thebibliography}{9}\footnotesize
\bibitem{bast2017} H.~Bast and C.~Korzen, A benchmark and evaluation for text extraction from PDF, \emph{JCDL} (2017).
\bibitem{parsebench} B.~Zhang \emph{et al.}, ParseBench, arXiv:2604.08538 (2026).
\bibitem{zhao2026} Y.~Zhao, S.~Pandey and P.~K.~Misra, What the author wrote, what the agent reads, arXiv:2610.30026 (2026).
\bibitem{liu2026} S.~Liu \emph{et al.}, Beyond the editing canvas, arXiv:2608.25880 (2026).
\bibitem{li2026} P.~Li, When file ingestion misleads language models, arXiv:2611.05806 (2026).
\end{thebibliography}

\clearpage
\onecolumngrid\raggedbottom
\setcounter{table}{0}\renewcommand{\thetable}{S\arabic{table}}
\section*{Supplement}
Everything below is generated from the same stored results as the main text (\texttt{make\_numbers.py}).

\subsection*{S1. File classes and coverage}
\begin{table}[H]
\caption{Constructed file classes (seeds make 20 files per class). PDF classes P2, P5, P8, P10, P12 and P13 exist in the generator and were not used for any number in this note.\label{tab:classes}}
\begin{ruledtabular}
\begin{tabular}{lp{0.8\textwidth}}
P1, P3, P4 & PDF: plain text pages (control); text pages with four blank pages among them; pages with rotation entries of 90, 180 or 270 degrees (the first file was defective and is withdrawn; P4c is the corrected file, Sec.~S5) \\
P6, P7, P9, P11 & PDF: image of the words with the same words as invisible text; image of the words only; form field values only; file cut at 50--70\% of its bytes \\
D1--D4 & DOCX: body paragraphs (control); tables; per-section headers and footers; table nested in a cell of a table \\
S1--S4 & PPTX: text boxes (control); table; grouped shapes; speaker notes \\
H1--H3 & HTML: paragraphs (control); table and list; navigation, aside and footer around an article \\
H4--H6 & HTML: text written by a script (the static text is core, the script-written words are scored separately); character entities; a block hidden by CSS \\
T1--T5 & Text: UTF-8 (control); UTF-16 with byte-order mark; Latin-1; cp1252; mixed UTF-8 and Latin-1 \\
\end{tabular}
\end{ruledtabular}
\end{table}
Coverage of the framework defaults (T = tested on all classes of the format listed above; U = not tested; the call is named in Sec.~II): LangChain: PDF T, DOCX T, HTML T, PPTX U, text U. LlamaIndex: all five T. Docling: PDF, DOCX, PPTX, HTML T; text U. Haystack: PDF, DOCX, HTML, text T; PPTX U (\texttt{PPTXToDocument} exists in 3.5.0; \texttt{MarkdownToDocument} also U). Unstructured and MarkItDown: all five T, default strategy only. Not tested: explicit strategy or OCR settings of Unstructured; the plain-settings PDF pipelines on the corrected rotated-page file.

\subsection*{S2. Tested cells without core loss and without rejection}
Table~S2 lists them.
\begin{table}[H]
\caption{Loader and classes that returned with token recall 1 on core content in every run (the controls are among them). The columns are as in Table~I.\label{tab:clean}}
\footnotesize
\begin{ruledtabular}
\begin{tabular}{llccc}
Default loader & File class & Lost & Raised & Recall \\
\colrule
LangChain \texttt{PyPDFLoader} & PDF, P1, P3, P6 & 0/60 & 0 & 1.00 \\
LangChain \texttt{Docx2txtLoader} & DOCX, all four classes & 0/80 & 0 & 1.00 \\
LlamaIndex \texttt{SimpleDirectoryReader} & TXT, UTF-8 (T1) & 0/20 & 0 & 1.00 \\
 & HTML, five classes (not H5) & 0/100 & 0 & 1.00 \\
 & PDF, P1, P3, P6 & 0/60 & 0 & 1.00 \\
 & DOCX and PPTX, all eight classes & 0/160 & 0 & 1.00 \\
Docling \texttt{DocumentConverter} & DOCX, D1, D2, D4 & 0/15 & 0 & 1.00 \\
 & PPTX, S1, S2, S3 & 0/15 & 0 & 1.00 \\
 & PDF, P1, P3 & 0/10 & 0 & 1.00 \\
 & HTML, all six classes & 0/30 & 0 & 1.00 \\
Haystack & \texttt{DOCXToDocument}, D1, D2 & 0/40 & 0 & 1.00 \\
 & \texttt{HTMLToDocument}, all six classes & 0/120 & 0 & 1.00 \\
 & \texttt{PyPDFToDocument}, P1, P3, P6 & 0/60 & 0 & 1.00 \\
 & \texttt{TextFileToDocument}, T1 & 0/20 & 0 & 1.00 \\
Unstructured & DOCX, all four classes & 0/80 & 0 & 1.00 \\
 & PPTX, S1, S2, S3 & 0/60 & 0 & 1.00 \\
 & HTML, all six classes & 0/120 & 0 & 1.00 \\
 & PDF, P1, P3, P6 & 0/60 & 0 & 1.00 \\
 & TXT, T1, T2 & 0/40 & 0 & 1.00 \\
MarkItDown & DOCX, D1, D2, D4 & 0/60 & 0 & 1.00 \\
 & PPTX, all four classes & 0/80 & 0 & 1.00 \\
 & HTML, all six classes & 0/120 & 0 & 1.00 \\
 & PDF, P1, P3, P6 & 0/60 & 0 & 1.00 \\
 & TXT, T1, T2 & 0/40 & 0 & 1.00 \\

\end{tabular}
\end{ruledtabular}
\end{table}

\subsection*{S3. Registered predictions against results: all \preregMiss\ misses}
The predictions (Haystack, Unstructured, MarkItDown; cell = loader $\times$ class) were committed with the plan before the run, from one throwaway development seed (20261400) and the framework documentation. There were \preregCells\ cells. They are the classes of Table~S1 that each loader was run on: the seven PDF classes P1, P3, P4, P6, P7, P9 and P11, the four DOCX classes, the six HTML classes and the five text classes for all three loaders (22 classes, \preregCellsHS\ cells for Haystack, which had no PPTX run), plus the four PPTX classes for Unstructured and MarkItDown (\preregCellsUN\ cells each), that is $22+26+26$. The cells for the original P4 (rotated pages) were registered for the defective file: Haystack no loss, Unstructured and MarkItDown silent loss. The count of \preregHit\ matches includes those three, which were scored on the original file; re-scoring the same three cells on the corrected file P4c (Table~S5) gives the same outcome (no loss for Haystack in 20 of 20 files; loss in 20 of 20 for Unstructured and 15 of 20 for MarkItDown, no exception), so the original tally of \preregHit/\preregCells\ is unchanged by P4c. The separate prediction made before the rerun (Amendment 5: no loss for any loader) is not part of that tally; it failed for Unstructured and MarkItDown. Of the six unreadable-content cells (P7 and P9 for each of the three loaders), five matched and one is a miss (Unstructured, P9). Of the \preregCells\ cells, \preregHit\ matched: \preregAHit\ registered silent core losses, \preregRaisedHit\ registered rejections, \preregBHit\ registered clean cells, and \preregUnreadHit\ registered unreadable cells (image-only and form PDFs). The seven that did not:
\begin{table}[H]
\caption{The seven wrong registered predictions (20 seeds each).\label{tab:prereg}}
\begin{ruledtabular}
\begin{tabular}{llll}
Loader & Class & Registered & Observed \\
\colrule
MarkItDown & TXT Latin-1 (T3) & recall 1 in 20 of 20 & lost tokens in 1 of 20, no exception \\
MarkItDown & TXT cp1252 (T4) & recall 1 in 20 of 20 & lost tokens in 3 of 20, no exception \\
MarkItDown & TXT mixed (T5) & recall 1 in 20 of 20 & lost tokens in 3 of 20, no exception \\
Unstructured & PDF form fields (P9) & form values lost & returned (unread in 0 of 20) \\
Unstructured & TXT Latin-1 (T3) & recall 1 in 20 of 20 & lost tokens in 1 of 20, no exception \\
Unstructured & TXT cp1252 (T4) & recall 1 in 20 of 20 & lost tokens in 3 of 20, no exception \\
Unstructured & TXT mixed (T5) & recall 1 in 20 of 20 & lost tokens in 3 of 20, no exception \\

\end{tabular}
\end{ruledtabular}
\end{table}
The text-file decoding of Unstructured and MarkItDown is not exact on short 8-bit files: in the same seven files for both loaders, \eightBitLostLo\ to \eightBitLostHi\% of the tokens were lost, about the share of words that hold a non-ASCII character; we did not examine the cause.

\subsection*{S4. Library calls}
\begin{table}[H]
\caption{Common library calls on constructed files: runs (of 20) that returned with core loss and with no exception and no new message on standard error, and pooled token recall; each as main run, hold-out run. The PyMuPDF truncated row is the case that has a flag: \texttt{Document.is\_repaired} was true for \pRepaired\ of \pFiles\ main-run files.\label{tab:lib}}
\footnotesize
\begin{ruledtabular}
\begin{tabular}{llcc}
Call & Class & No exception or new stderr message & Recall \\
\colrule
python-docx paras. & tables (D2) & 20, 20 & 0.71, 0.72 \\
 & header, footer (D3) & 20, 20 & 0.82, 0.81 \\
 & nested table (D4) & 20, 20 & 0.76, 0.78 \\
mammoth raw text & header, footer (D3) & 20, 20 & 0.82, 0.81 \\
python-pptx frames & tables (S2) & 20, 20 & 0.76, 0.80 \\
 & grouped shapes (S3) & 20, 20 & 0.67, 0.68 \\
 & speaker notes (S4) & 20, 20 & 0.58, 0.58 \\
PyMuPDF \texttt{get\_text} & truncated: damaged-input handling (P11)$^{a}$ & 20, 19 & 0.53, 0.52 \\

\end{tabular}
\end{ruledtabular}
\raggedright\footnotesize $^{a}$ Repair flag set in \pRepaired\ of \pFiles\ main-run files; in the hold-out run one file raised.

\medskip\centering\emph{Framework defaults on truncated PDFs (P11): damaged-input handling, recall against the intact original; columns as in Table~I.}
\footnotesize
\begin{ruledtabular}
\begin{tabular}{llccc}
Default loader & File class & Lost, no exception & Raised & Recall \\
\colrule
LangChain \texttt{PyPDFLoader} & PDF, truncated (P11) & 0/20 & 20 & -- \\
LlamaIndex \texttt{SimpleDirectoryReader} & PDF, truncated (P11)$^{c}$ & 20/20 & 0 & 0.00 \\
Docling \texttt{DocumentConverter} & PDF, truncated (P11)$^{d}$ & 5/5 & 0 & 0.49 \\
Haystack \texttt{PyPDFToDocument} & PDF, truncated (P11)$^{b}$ & 20/20 & 0 & 0.00 \\
Unstructured \texttt{partition} & PDF, truncated (P11) & 0/20 & 20 & -- \\
MarkItDown \texttt{convert} & PDF, truncated (P11) & 0/20 & 20 & -- \\

\end{tabular}
\end{ruledtabular}
\raggedright\footnotesize $^{b}$ Haystack skips the file, writes a warning to the log and returns an empty \texttt{documents} list (zero documents, checked); the call does not raise. $^{c}$ No exception and an empty document list; in a one-file check the reader printed ``Failed to load file \ldots\ Skipping'' to standard output. $^{d}$ Checked afterwards (5 of 5 files): the result status was \texttt{PARTIAL\_SUCCESS} with 3 to 6 entries in \texttt{errors}, so a caller that inspects the status is told.
\end{table}

\subsection*{S5. Rotated pages: the clip check and the corrected rerun (Amendment 5)}
\emph{The defect.} Over the 20 framework seeds, PyMuPDF \texttt{get\_text()} returned \pFourDefault\ of the \pFourWords\ known words of the first P4 file by default and \pFourNoClip\ with \texttt{clip=pymupdf.INFINITE\_RECT()}; all \pFourExtra\ words returned only without the clip lay entirely outside the page rectangle, and CropBox equalled MediaBox on all \pFourPages\ pages. Pages with a rotation entry of 90 or 270 (\pFourNinety\ pages) had MediaBox $792\times612$ with text drawn up to 740 points above the lower edge. No extraction result from that file appears in the results tables (the diagnostic counts above are the only numbers taken from it); the three registered P4 cells in Sec.~S3 were scored on it first and are retained in that tally (script \texttt{stage3/p4\_clip\_check.py}).

\emph{The rerun.} \texttt{stage2/gen2\_p4fix.py} differs from the generator in one constant (text starts 560 points up on rotated pages); same seeds gave the same words (checked against the stored manifests). The plan and the prediction (no loss for any loader) were committed before the run; the prediction failed for Unstructured (loss in 20 of 20 files) and MarkItDown (15 of 20) and held for the other loaders and PyMuPDF. Default calls as in Sec.~II, 20 seeds (Docling 5), \texttt{stage3/p4fix\_run.py}.
\begin{table}[H]
\caption{Rotated pages, corrected file (P4c). Columns as in Table~I. PyMuPDF with its default \texttt{get\_text()}.\label{tab:pfour}}
\footnotesize
\begin{ruledtabular}
\begin{tabular}{llccc}
Default loader & File class & Lost & Raised & Recall \\
\colrule
LangChain \texttt{PyPDFLoader} & PDF, rotated pages, corrected file (P4c) & 0/20 & 0 & 1.00 \\
LlamaIndex & PDF, rotated pages, corrected file (P4c) & 0/20 & 0 & 1.00 \\
Docling & PDF, rotated pages, corrected file (P4c) & 0/5 & 0 & 1.00 \\
Haystack \texttt{PyPDFToDocument} & PDF, rotated pages, corrected file (P4c) & 0/20 & 0 & 1.00 \\
Unstructured & PDF, rotated pages, corrected file (P4c) & 20/20 & 0 & 0.39 \\
MarkItDown & PDF, rotated pages, corrected file (P4c) & 15/20 & 0 & 0.61 \\
PyMuPDF \texttt{get\_text} & PDF, rotated pages, corrected file (P4c) & 0/20 & 0 & 1.00 \\

\end{tabular}
\end{ruledtabular}
\end{table}
For Unstructured and MarkItDown the text is inside the page (PyMuPDF with its page clip read all of it). As a standalone diagnostic (the wrapper's strategy was not recorded): on one page with a rotation entry of 90, \texttt{pdfminer.six} with its default \texttt{LAParams} returned word recall 0.00 and with \texttt{detect\_vertical=True} 1.00 (\texttt{stage3/p4fix\_pdfminer\_check.py}); Package metadata lists \texttt{pdfminer.six} as a dependency of both packages. The one-page check was run directly with \texttt{pdfminer.six}, and the Unstructured strategy used on P4c was not recorded.

\emph{MarkItDown results by rotation sequence.} The generator draws one rotation sequence per seed from four (all pages 90; 270 and 0; 180 and 90; 90, 180, 270 and 0 in turn). Regenerating the 20 files from the stored generator and seeds (the manifest for seed 20261401 is identical to the stored one, and the 20 files hold 18,829 known words, the denominator of the pooled recall in Table~S5) shows that the five MarkItDown files that returned every word (seeds 20261402, 20261403, 20261404, 20261409 and 20261419) are exactly the five files in which every page is rotated by 90 degrees; Unstructured returned no word from those five. The other 15 files contain pages rotated by 180 or 270 degrees, and MarkItDown lost words in all of them. MarkItDown recovered all tokens from the five all-90 files, while Unstructured recovered none. The stored results do not localize MarkItDown's omissions within the mixed-rotation files, so we make no claim about 180 or 270 degree pages for MarkItDown; the one-page pdfminer check does not explain its pattern and we did not isolate the cause.


\subsection*{S6. Real documents}
Public files: 59 arXiv PDFs (January 2026, six categories), 50 PDFs, 20 HTML and 20 text files from a GovDocs1 subset, and 40 Wikipedia pages, with source, date and checksum recorded. A file enters a comparison only if its reference has at least 50 tokens. For PDFs the reference is pypdfium2: 103 of the 109 PDFs qualified (the other 6 gave fewer than 50 tokens, mostly scans). LangChain's \texttt{PyPDFLoader} and LlamaIndex's reader each returned fewer than 90\% of the reference tokens in 16 of 103 files (Wilson 95\%: 0.10 to 0.24) and fewer than 50\% in 1. Docling's PDF run was stopped for running time after the first 30 PDFs; 18 of 30 had under 90\% of the reference tokens. Docling produces a different text model (reading order, hyphenation, tables), so this is divergence between tools and we do not call it loss. For web pages the reference is the separator call; 58 of 60 pages qualified (one page returned no text from any reader, one had 27 tokens). The 20 GovDocs1 text files were read with LlamaIndex's reader only, against the fullest of the registered text idioms (strict decoding with byte-order mark): all 20 qualified, and none of the 20 files fell below 90\% overlap with the reference (0 of 20; Wilson 95\% upper bound 0.16).

\subsection*{S7. A general check that we tried and withdraw}
A run-time check that compares a loader's tokens with those of a second, independent reader and withholds the output on disagreement was tested for PDF. On the 109 real PDFs it withheld or flagged all 109, mostly because the second reader failed on real fonts (76 files). It also accepted the incomplete page in 540 of 540 constructed runs where a page mixed readable text with text drawn as an image or as outlines, because both readers missed the same visible content. In the tested PDF setup their agreement therefore did not establish complete extraction. We keep only the measurement.

\subsection*{S8. Versions and run details}
\begin{table}[H]
\caption{Versions, read from the installed packages. Python 3.12.13, macOS, CPU only. The framework environment and the library environment are separate; pypdf differs between them.\label{tab:versions}}
\begin{ruledtabular}
\begin{tabular}{lp{0.82\textwidth}}
Frameworks & langchain-community 0.4.2, langchain-core 1.8.4; llama-index-core 0.14.26, llama-index-readers-file 0.8.0; docling 2.152.0, docling-core 2.108.0, docling-parse 7.27.1; with pypdf 6.21.4, beautifulsoup4 4.15.0, docx2txt 0.9; haystack-ai 3.5.0, unstructured 0.27.38, markitdown 0.1.9 (own environment: trafilatura 2.4.0, pdfminer.six 20260107, python-docx 1.2.0, python-pptx 1.0.2, lxml 6.1.3) \\
Libraries & PyMuPDF 2.1 (MuPDF 1.28.6), pypdf 7.1.1, pypdfium2 5.15.2, pikepdf 10.18.2, python-docx 1.2.0, python-pptx 1.0.2, mammoth 1.13.0, docx2txt 0.9, markitdown 0.1.9, beautifulsoup4 4.15.0, lxml 6.1.3 \\
Generator & reportlab 5.1.0, python-docx 1.2.0, python-pptx 1.0.2 \\
\end{tabular}
\end{ruledtabular}
\end{table}
Framework seeds were 20261401 to 20261420. Unstructured with \texttt{auto} ran OCR on the image-only pages, where it returned 545 of 8,905 known words (6\%) from the 20 files; the strategy chosen on text-layer files was not recorded.

\subsection*{S9. Docling afterwards: Markdown export and result status}
Post hoc, not registered (\texttt{stage3/docling\_md\_check.py}, \texttt{stage3/status\_checks.py}), seeds 20261401 to 20261405. Pooled core recall of \texttt{export\_to\_markdown()} equalled that of \texttt{export\_to\_text()} on D3 (0.83, 537 of 647) and S4 (0.57, 768 of 1350). \texttt{ConversionResult.status} was \texttt{SUCCESS} with no errors for every class except P11, where it was \texttt{PARTIAL\_SUCCESS} with 3 to 6 errors in 5 of 5 files. Haystack: the skipped inputs (T2 to T5 and P11) returned an empty \texttt{documents} list (zero documents; \texttt{TextFileToDocument} on the UTF-8 control returned one).

\subsection*{S10. Alternative settings: counts by call}
Runs attempted on the constructed files, as complete / incomplete / rejected / no text layer (image-only and form PDFs have no core tokens). PDF counts are over the classes P1, P3, P6, P7, P9 and P11 (P4 withdrawn; P4c was not run with these calls).
\begin{table}[H]
\caption{Alternative settings, counts by the call used. ``Rejected'' are runs that raised.\label{tab:alt}}
\footnotesize
\begin{ruledtabular}
\begin{tabular}{lccccc}
Call & Complete & Incomplete & Rejected & No text layer & Runs \\
\colrule
BeautifulSoup \texttt{get\_text(" ")}, HTML & \pHtmlComplete & \pHtmlIncomplete & \pHtmlRejected & \pHtmlNocore & \pHtmlRuns \\
python-docx walk with tables, sections, headers and footers, DOCX & \pDocxComplete & \pDocxIncomplete & \pDocxRejected & \pDocxNocore & \pDocxRuns \\
python-pptx walk with groups, tables and notes, PPTX & \pPptxComplete & \pPptxIncomplete & \pPptxRejected & \pPptxNocore & \pPptxRuns \\
strict decoding with byte-order mark, text (rejections are T3--T5) & \pTxtComplete & \pTxtIncomplete & \pTxtRejected & \pTxtNocore & \pTxtRuns \\
pypdfium2 with end-of-file and \texttt{pikepdf} syntax checks, PDF & \pPdfBComplete & \pPdfBIncomplete & \pPdfBRejected & \pPdfBNocore & \pPdfBRuns \\
PyMuPDF with the repair check, PDF & \pPdfAComplete & \pPdfAIncomplete & \pPdfARejected & \pPdfANocore & \pPdfARuns \\
\end{tabular}
\end{ruledtabular}
\end{table}
For both PDF pipelines the 20 rejections are the truncated files, and the 40 runs with no text layer are the image-only and form classes; image-only pages came back empty in \cUnreadPdf\ of 20 files for both, and the words written by a script in H4 were not returned.


\end{document}
